\documentclass{amsart}
\usepackage[a4paper,left=3cm,right=3cm,top=3cm,bottom=4cm]{geometry}
\usepackage{amsmath,amssymb,amsthm}
\usepackage{microtype}
\usepackage[colorlinks=true,linkcolor=blue,citecolor=blue,urlcolor=blue]{hyperref}

\newtheorem{theorem}{Theorem}
\newtheorem{lemma}{Lemma}[section]
\newtheorem{remark}{Remark}[section]
\DeclareMathOperator{\Aut}{Aut}
\DeclareMathOperator{\Sym}{Sym}
\DeclareMathOperator{\Hnm}{H}
\DeclareMathOperator{\Jac}{Jac}
\DeclareMathOperator{\rk}{rk}
\setlength{\parindent}{15pt}
\setlength{\parskip}{0pt plus 1pt}
\allowdisplaybreaks[1]

\title[The tangent Gauss map]{A counterexample to generic injectivity of the tangent Gauss map}
\author{Xiaoxiang Zhou}
\address{Institut f\"ur Mathematik, Universität Hamburg,
20146 Hamburg, Germany}
\email{xiaoxiang.zhou@uni-hamburg.de}
\date{\today}

\begin{document}
\maketitle

We work over $\mathbb{C}$. For a curve $C\subset A$ in an abelian variety,
the tangent Gauss map sends a smooth point to its tangent line, translated
to the origin:
\[
\phi_C:C^{\mathrm{sm}}\longrightarrow\mathbb{P}(T_0A),\qquad
p\longmapsto d(t_{-p})_p(T_pC).
\]
We call $C$ nondegenerate if it is not contained in a translate of a
proper abelian subvariety.

\begin{theorem}\label{thm:weak}
There exists a nondegenerate curve in an abelian threefold whose tangent
Gauss map has generic degree $3$.
\end{theorem}

\section{Strategy}

Let $C$ be a smooth non-hyperelliptic curve of genus $4$, and let
$\mathcal{L}$ be a line bundle of degree $3$ with $h^0(\mathcal{L})=2$.
It is base-point-free, since a base point would give a degree-two map to
$\mathbb{P}^1$.
Thus $|\mathcal{L}|:C\longrightarrow\mathbb{P}^1$ has degree $3$.
The linear system
$\left(\mathcal{L}^{\otimes2},\Sym^2\Hnm^0(\mathcal{L})\right)$
defines the map
\[
\phi:C\overset{|\mathcal{L}|}{\longrightarrow}\mathbb{P}^1
\overset{\nu_2}{\longrightarrow}\mathbb{P}^2,
\]
where $\nu_2$ is the quadratic Veronese embedding. In particular, $\phi$
has generic degree $3$ onto its image and is not generically injective.

To obtain a counterexample by this construction, we seek a pair
$(C,\mathcal{L})$ for which this map is a tangent Gauss map. It suffices
to impose two further conditions:
\begin{enumerate}
\item\label{cond:spin} $\mathcal{L}^{\otimes2}\cong\omega_C$;
\item\label{cond:hodge} $\Sym^2\Hnm^0(\mathcal{L})\subset\Hnm^0(\omega_C)$ is the
$(1,0)$-part of a rational Hodge substructure of
$\Hnm^1(C;\mathbb{Q})$.
\end{enumerate}
The second condition allows these three canonical forms to be pulled
back from invariant forms on an abelian threefold; we make this passage
in Section~\ref{subsec:image}.

If we temporarily drop nonhyperellipticity, both conditions can be met.
Take a smooth hyperelliptic curve $C_0$ of genus $4$ admitting a
tangential cover $\pi_0:(C_0,p_0)\longrightarrow(E_0,0)$ at a Weierstrass
point $p_0$. Such examples are supplied by the Lam\'e spectral curves
\cite[Table~3 and Theorem~L]{maier} and \cite[Proposition~3.1]{eilbeck}.
Tangentiality means that $\pi_0^*E_0\subset\Jac(C_0)$ is tangent to the
Abel--Jacobi image at $p_0$. Put
$\mathcal{L}_0=\mathcal{O}_{C_0}(3p_0)$. Then
\[
\mathcal{L}_0^{\otimes2}\cong\omega_{C_0},\qquad
\Sym^2\Hnm^0(\mathcal{L}_0)
=\Hnm^0\left(\omega_{C_0}(-2p_0)\right).
\]
Indeed, a hyperelliptic coordinate $x$ with a double pole at $p_0$ gives
the basis $1,x$ of $\Hnm^0(\mathcal{L}_0)$, with products $1,x,x^2$.
Tangentiality gives the Hodge condition, as explained in
Lemma~\ref{lem:hyperelliptic}. This seed satisfies the two conditions,
but $\mathcal{L}_0$ has a base point at $p_0$. We therefore deform the
pair while preserving both conditions and leaving the hyperelliptic
locus.

\section{A dimension count for deformations}

We first preserve the spin condition and the two sections, then impose
compatibility of two different deformations of the cohomology subspaces.
A dimension comparison will allow us to remove the hyperelliptic locus.

\subsection{The deformation space and the spin condition}\label{subsec:spin}

Let $B$ be a small simply connected marked deformation space of the spin
curve $(C_0,\mathcal{L}_0)$. Thus
$\mathcal{L}_t^{\otimes2}\cong\omega_{C_t}$ throughout $B$, and
$\dim B=3g-3=9$. To retain two sections, restrict to the theta-null locus
\[
S=\left\{t\in B\;\middle|\;h^0(\mathcal{L}_t)\geqslant2\right\}
\]
locally defined by one theta constant. Hence $\dim S\geqslant8$.
After shrinking $B$, upper semicontinuity gives
$h^0(\mathcal{L}_t)=2$ on $S$.

\subsection{Compatibility of the two deformations}\label{subsec:compatibility}

At the seed, set
\[
\begin{aligned}
V_0&:=\left(\pi_0^*\Hnm^1(E_0;\mathbb{Q})\right)^\perp
\subset\Hnm^1(C_0;\mathbb{Q}),\\
W_0&:=\Sym^2\Hnm^0(\mathcal{L}_0)\subset\Hnm^0(\omega_{C_0}),
\end{aligned}
\]
where $\perp$ is taken with respect to the cup-product pairing.
Tangentiality gives $V_{0,\mathbb{C}}=W_0\oplus\overline{W}_0$.
Over $S$, we deform these subspaces in two different ways:
\[
\begin{aligned}
V_0&\longmapsto V_t\subset\Hnm^1(C_t;\mathbb{Q})
&&\text{by Gauss--Manin parallel transport},\\
W_0&\longmapsto W_t:=\Sym^2\Hnm^0(\mathcal{L}_t)
&&\text{by deforming the line bundle}.
\end{aligned}
\]
Thus $V_t$ is a flat rational subspace of dimension $6$, whereas
$W_t\subset\Hnm^0(\omega_{C_t})$ is a holomorphic subbundle of rank $3$.
These constructions need not remain compatible: $V_t$ may cease to be
a Hodge substructure, and $W_t$ need not lie in its complexification.
Their compatibility is precisely the condition we impose:
\[
\begin{aligned}
&\begin{cases}
V_t\text{ remains a Hodge substructure},\\
V_{t,\mathbb{C}}\cap\Hnm^0(\omega_{C_t})=W_t
\end{cases}\\
\Longleftrightarrow\;& V_{t,\mathbb{C}}=W_t\oplus\overline{W}_t\\
\Longleftrightarrow\;& W_t\subseteq V_{t,\mathbb{C}}\\
\Longleftrightarrow\;&
\rk\left[W_t\longrightarrow\Hnm^1(C_t;\mathbb{C})/V_{t,\mathbb{C}}\right]=0.
\end{aligned}
\]
The middle equivalence follows from rationality of $V_t$ and Hodge type,
since $\dim W_t=3$ and $\dim V_t=6$. The last condition is the vanishing
of a bundle map of ranks $3$ and $2$, so it imposes at most six equations
on $S$. Consequently, every irreducible component $Z$ of its zero locus
through the seed satisfies
\[
\dim Z\geqslant9-1-6=2.
\]
Fix one such component $Z$.

\subsection{Removing the hyperelliptic locus}\label{subsec:hyperelliptic}

To find a nonhyperelliptic point, it remains to bound
$\dim\left(Z\cap\mathcal{H}_4\right)$, where $\mathcal{H}_4\subset B$
is the hyperelliptic locus. The following lemma reduces this to counting
tangential covers.

\begin{lemma}\label{lem:hyperelliptic}
Let $C$ be a smooth hyperelliptic curve of genus $4$, and let
$\mathcal{L}^{\otimes2}\cong\omega_C$ with $h^0(\mathcal{L})=2$.
\begin{enumerate}
\item There is a Weierstrass point $p$ such that $\mathcal{L}\cong\mathcal{O}_C(3p)$.
\item For this point $p$, the subspace
$\Sym^2\Hnm^0(\mathcal{L})=\Hnm^0\left(\omega_C(-2p)\right)$ is the $(1,0)$-part of a rational
Hodge substructure if and only if $C$ admits a tangential cover at $p$.
\end{enumerate}
\end{lemma}

\begin{proof}\
\begin{enumerate}
\item Fix a Weierstrass point $w$. If $\mathcal{L}$ were base-point-free, its
degree-three map and the degree-two map $|\mathcal{O}_C(2w)|$ would give
a birational map to a curve of bidegree $(2,3)$ in
$\mathbb{P}^1\times\mathbb{P}^1$, of arithmetic genus $2$. Thus
$\mathcal{L}$ has a base point $p$, and uniqueness of the degree-two
linear system gives $\mathcal{L}(-p)\cong\mathcal{O}_C(2w)$.
Using $\omega_C\cong\mathcal{O}_C(6w)$, the spin condition gives
$\mathcal{O}_C(2p)\cong\mathcal{O}_C(2w)$. Hence $p$ is Weierstrass and
$\mathcal{L}\cong\mathcal{O}_C(3p)$. 

\item Under $T_0^*\Jac(C)=\Hnm^0(\omega_C)$, the space
$\Hnm^0\left(\omega_C(-2p)\right)=\Hnm^0\left(\omega_C(-p)\right)$
is the annihilator of the tangent line at $p$. Its being the $(1,0)$-part
of a rational Hodge substructure is therefore equivalent to this line
being the tangent space of an elliptic subvariety of $\Jac(C)$.
Dualizing the inclusion and composing with the Abel--Jacobi map gives
a tangential cover; the converse follows from the same annihilator
description.
\end{enumerate}
\end{proof}

Consequently, the hyperelliptic points of $Z$ give tangential covers
$\pi_t:(C_t,p_t)\longrightarrow(E_t,0)$. The complementary Hodge
substructure $V_t^\perp$ has rank $2$ and determines $E_t$ up to isogeny.
Choosing its saturated integral lattice fixes the pullback map on
integral cohomology under local markings. The cup-product pairing then
shows that $\deg\pi_t$ is locally constant.

\begin{remark}
For a fixed elliptic curve and degree there are only finitely many
hyperelliptic tangential covers
\cite[Proposition~2.5, 3.; Section~2, Remark~1, 3.--4.]{treibich2025}.
Indeed, the relation $2n+1=\sum_{i=0}^3\mu_i^2+4d$ leaves finitely many
types for a fixed degree $n$, and each corresponding Severi set is
finite. 
\end{remark}

Since elliptic curves have one modulus and the choices of Weierestrass points
$p_t$ are finite,
\[
\dim\left(Z\cap\mathcal{H}_4\right)\leqslant1<2\leqslant\dim Z.
\]
We can therefore choose a nonhyperelliptic point of $Z$.

\subsection{Passing to a curve in an abelian threefold}\label{subsec:image}

For this point, the rational Hodge substructure $V$ gives a quotient
$q:\Jac(C)\longrightarrow A$ with connected kernel and
$q^*\Hnm^0(\Omega_A^1)=\Sym^2\Hnm^0(\mathcal{L})$.
Compose the Abel--Jacobi map based at $b\in C$ with $q$:
$f=q\circ j_b:C\longrightarrow A$. The pulled-back invariant forms
are the products of sections of $\mathcal{L}$, so $f$ is immersive and
its tangent map is $\phi$. Its image contains the origin and generates
$A$, hence is nondegenerate. If $f$ had degree $e>1$ onto the
normalization of its image, then $e\mid3$, since $\phi$ factors through
that normalization. Thus $e=3$, and the normalization would map with
degree $1$ to the conic $\nu_2(\mathbb{P}^1)$. This is impossible in an
abelian variety. Therefore $C$ is the normalization of $f(C)$ and the
tangent Gauss map of $f(C)$ has generic degree $3$, proving
Theorem~\ref{thm:weak}.

\section{Simplicity and symmetries}

\begin{theorem}\label{thm:strong}
The abelian threefold can be chosen simple. The resulting curve is
preserved by no nonidentity automorphism of $A$ with scalar differential.
\end{theorem}

First choose the hyperelliptic seed so that
$A_0:=\Jac(C_0)/\pi_0^*E_0$ is simple; the existence of such a seed is
checked below. For every nonzero proper rational
subspace $R\subset V_0$ on which the polarization is nondegenerate, let
$e_R$ be the orthogonal projector. The condition
\[
e_R(W_t)\subset W_t
\]
defines a closed analytic subset of $Z$. It is proper because it misses
the seed, and it holds precisely when $R$ is a Hodge substructure.
There are only countably many $R$, so Baire's theorem allows us to choose
a point outside these subsets and $\mathcal{H}_4$. Its Hodge structure
$V_t$ is simple, hence the corresponding abelian threefold $A$ is simple.

Now suppose $\tau\in\Aut(A)$ preserves $f(C)$ and has scalar
differential, using translations to identify tangent spaces. It lifts
to an automorphism $\sigma$ of $C$ with $\phi\circ\sigma=\phi$.
If $\sigma$ is the identity, then $\tau$ fixes a generating curve
pointwise and is the identity. Otherwise, the order of $\sigma$ divides
$\deg\phi=3$, so it is $3$ and
$C/\langle\sigma\rangle\cong\mathbb{P}^1$. Riemann--Hurwitz gives a
fixed point of $\sigma$. Taking its image as the origin makes $\tau$
a group automorphism $u$ of order $3$. Its scalar differential is a
primitive cube root of unity $\lambda$.
Then $\mathbb{Q}(\lambda)$ acts on $\Hnm_1(A;\mathbb{Q})$, and the
complex structure on real homology is
$\left(u_*-\operatorname{Re}(\lambda)\right)/\operatorname{Im}(\lambda)$.
Every one-dimensional $\mathbb{Q}(\lambda)$-subspace therefore has a
complex-invariant realification and defines an elliptic subvariety of
$A$, contradicting simplicity. This proves the symmetry assertion.

\medskip
\noindent\emph{Existence of a seed with simple quotient.}
Take $E_0:y^2=4x^3-4x-1$, let $e_1,e_2,e_3$ be the roots of
$4e^3-4e-1$, and take the smooth projective model of
\[
C_0:\quad \nu^2=(B^3-208B+560)
\prod_{i=1}^3\left(B^2+10e_iB-35e_i^2-28\right).
\]
This is the genus-four Lam\'e spectral curve, with a degree-ten cover
of $E_0$ tangential at infinity
\cite[Table~3 and Theorem~L]{maier} and \cite[Proposition~3.1 and the discussion
after (17)]{eilbeck}.
The cover is defined over $\mathbb{Q}$, and $E_0$ has discriminant $37$.
Modulo $11$, the polynomial on the right is the square-free polynomial
$B^9+B^4+6B^3+9B^2+7B+5$; hence $C_0$ has genus $4$ and good reduction.
Direct counting gives
\[
\left(\#C_0(\mathbb{F}_{11^r})\right)_{r=1}^4
=(10,120,1312,15312),\qquad \#E_0(\mathbb{F}_{11})=7.
\]
Newton's identities and division by the elliptic factor $T^2-5T+11$
give the Frobenius polynomial of $A_0=\Jac(C_0)/\pi_0^*E_0$:
\[
h(T)=T^6+3T^5+5T^4-14T^3+55T^2+363T+1331.
\]
It is irreducible modulo $3$, and $11\nmid14$, so this reduction is
simple and ordinary. An ordinary simple abelian threefold over a
finite field of odd cardinality is absolutely simple unless its
Frobenius polynomial has the form $T^6+cT^3+q^3$
\cite[Proposition~3.4]{arias}. Thus this reduction is absolutely simple,
and so is $A_0$ in characteristic zero, since a proper abelian
subvariety would specialize to one after a finite extension.
% The existing verify_tangent_seed.py reproduces the reduction, point
% counts, Frobenius factorization, and irreducibility checks.

\begin{thebibliography}{9}
\raggedright
\bibitem{arias}
S.~Arias-de-Reyna, C.~Armana, V.~Karemaker, M.~Rebolledo,
L.~Thomas, and N.~Vila,
\emph{Large Galois images for Jacobian varieties of genus 3 curves},
\href{https://arxiv.org/abs/1507.05913}{arXiv:1507.05913}, 2016.

\bibitem{eilbeck}
J.~C.~Eilbeck, V.~Z.~Enolskii, and E.~Previato,
\emph{Varieties of elliptic solitons},
J.\ Phys.\ A: Math.\ Gen.\ \textbf{34} (2001), 2215--2227,
\href{https://doi.org/10.1088/0305-4470/34/11/314}{doi:10.1088/0305-4470/34/11/314}.

\bibitem{maier}
R.~S.~Maier,
\emph{Lam\'e polynomials, hyperelliptic reductions and Lam\'e band structure},
\href{https://arxiv.org/abs/math-ph/0309005v4}{arXiv:math-ph/0309005v4}, 2004.

\bibitem{treibich2025}
A.~Treibich,
\emph{Hyperelliptic tangential covers and even elliptic finite-gap
potentials, back and forth},
\href{https://arxiv.org/abs/2501.16483v1}{arXiv:2501.16483v1}, 2025.
\end{thebibliography}
\end{document}
